\section{Discussão}\label{sec:discussao}

Os resultados indicam que a qualidade das explicações GradientSHAP não acompanha o desempenho do modelo de forma única: depende da propriedade medida. As métricas que envolvem a decisão do modelo entre classes acompanharam o F1 nas pontas, enquanto as de fidelidade e de robustez seguiram outro padrão, ligado ao tipo de evidência de cada classe.

\subsection{Fidelidade e robustez seguem a evidência, e não o F1}
Os valores de fidelidade são compatíveis com a literatura. Em quatro bases de ECG, a \FC\ do DeepSHAP, outra aproximação dos valores SHAP, ficou entre 0,41 e 0,63, a maior entre os onze métodos avaliados~\cite{suh2024}; e, na mesma base deste estudo, a relevância atribuída pelos métodos funcionou mais como uma ordem de importância do que como uma contribuição proporcional~\cite{turbe2023}, o que torna esperadas correlações moderadas como as obtidas (0,60 e 0,61).

A ordem das classes na fidelidade concorda com o viés de amplitude descrito para os métodos que multiplicam o gradiente pela entrada ou subtraem uma referência, que deslocam a relevância para os componentes de maior amplitude do ECG~\cite{gumpfer2026}. No mesmo trabalho, esses métodos, entre eles o GradientSHAP, reforçaram as regiões de grande amplitude e falharam em cobrir o segmento ST na isquemia~\cite{gumpfer2026}. No presente estudo, as classes reconhecidas pela forma do QRS, o componente de maior amplitude do traçado, ficaram no topo da fidelidade, e as classes cuja evidência é um desvio pequeno (STD, onda P na CAP) ou uma propriedade da sequência de batimentos (a irregularidade do ritmo na AF) ficaram embaixo. Em dados sintéticos com evidência conhecida, variações simples de amplitude ou de duração do padrão entre as classes bastaram para produzir efeitos dependentes da classe, e as métricas de perturbação e as baseadas na evidência conhecida frequentemente avaliaram as classes de forma contraditória~\cite{baer2025b}. A diferença de fidelidade entre as classes não deve, portanto, ser lida por si como diferença de qualidade das explicações; Baer et al.\ mostram que a eficácia das perturbações varia muito entre as classes e recomendam complementar as médias agregadas com análises estratificadas por classe~\cite{baer2025a}. A AF ilustra o caso: retirar uma região não retira a irregularidade dos intervalos RR, e a queda da saída deixa de acompanhar a atribuição.

Na robustez, a principal fonte de variação das explicações foi o próprio método. As sensibilidades ficaram cerca de quatro vezes abaixo do piso, e a \MS\ caiu de 0,072 para 0,045 ao passar de 50 para 200 sorteios, o que mostra que parte da sensibilidade medida com poucos sorteios era erro de estimativa. Esse comportamento é coerente com a comparação de Mercier et al., em séries temporais, em que os métodos baseados em SHAP estiveram entre os de maior sensibilidade e os métodos de gradiente foram descritos como ruidosos e fracos nesse critério~\cite{mercier2022}, e mostra por que a leitura em relação ao piso é necessária: sem ela, a variação do método seria interpretada como falta de robustez da explicação. As classes menos robustas (CAP, STD e AF) foram também as de piso mais alto, isto é, aquelas em que o próprio método é menos estável, de novo as classes cuja evidência é pequena ou distribuída no tempo.

\subsection{Consistência e especificidade acompanham o F1}
A \CS\ foi a métrica cuja ordem mais se aproximou da ordem do F1. Duas propriedades da métrica contribuem para isso. Ela usa as predições do próprio modelo~\cite{dasgupta2022}, de modo que uma classe reconhecida por um padrão estável tende a produzir explicações repetidas e, ao mesmo tempo, um F1 alto. E ela depende do número de mapas de cada classe: BRD, com 172 mapas, teve o maior valor ao acaso (0,23), e STE, com 14, um valor ao acaso de 0,02, o que exige ler cada classe junto com o seu valor ao acaso.

No \RL, todos os pares ficaram acima do piso, de modo que os mapas não são o mesmo mapa para qualquer classe, falha que motivou a proposta da métrica~\cite{sixt2020}. Em EEG simulado, métodos da família gradiente × entrada também foram sensíveis à classe e à randomização do modelo~\cite{ravindran2023}. O fato de STE, a classe de menor F1, ter as explicações menos específicas é coerente com um modelo que separa pior essa classe das demais: se a evidência usada para STE se parece com a usada para outras classes, os mapas também se parecem.

\subsection{Complexidade e dependência do modelo}
A relação da complexidade com o F1 foi apenas parcial, e a complexidade foi a categoria de métricas menos confiável numa meta-avaliação das métricas~\cite{hedstrom2023metaquantus}, que recomenda não estimar a qualidade das explicações por uma única categoria. Por isso, ela foi lida em conjunto com as demais.

O \EM\ não foi satisfeito: o mapa do modelo aleatório ficou mais concentrado, e não mais espalhado. Isso é esperado num método que multiplica o gradiente por (ECG $-$ referência), Eq.~\eqref{eq:gshap}: sem o que o modelo aprendeu, o segundo fator continua presente e é maior nos picos do QRS, de modo que o mapa reproduz a forma do próprio ECG. Para métodos do tipo gradiente × entrada, a randomização dos pesos altera o sinal das atribuições, mas a estrutura da entrada permanece nos mapas~\cite{adebayo2018}; no ECG, o mesmo viés coloca a relevância no pico R mesmo em tarefas sobre as ondas P e T~\cite{wagner2024}. A entropia não capta essa diferença, mas as distâncias captam: o mapa do modelo aleatório ficou a 8,7 vezes o piso do mapa treinado, o que concorda com a perda rápida de correlação entre os mapas à medida que os pesos são randomizados, observada em ResNets para ECG~\cite{arends2026}. O mapa depende, portanto, do modelo, e o \EM\ foi tomado como teste de sanidade, e não como métrica de comparação entre classes.

\subsection{Resposta à pergunta do estudo}
Os resultados respondem à pergunta de forma condicional. A qualidade das explicações acompanha o desempenho do modelo nas propriedades que dependem da forma como o modelo separa as classes, a consistência entre explicação e predição e a especificidade da explicação para a classe, a mesma capacidade que o F1 mede. Ela não acompanha o desempenho na fidelidade e na robustez, que dependem da amplitude e da extensão da evidência de cada classe e, no GradientSHAP, do viés do método para os componentes de grande amplitude do ECG. Para o uso clínico das explicações, isso significa que o F1 de uma classe não indica, sozinho, se as suas explicações são fiéis ao modelo.

\subsection{Limitações}
Retirar regiões cria ECGs diferentes dos vistos no treino, e a queda da saída pode refletir, em parte, essa diferença e não apenas a perda de informação~\cite{ismail2020}. Perturbar trechos contíguos introduz menos ruído de alta frequência que perturbar pontos espalhados~\cite{cui2023}, e as duas referências testadas deram o mesmo resultado, mas o efeito não foi eliminado. Foi avaliado um único modelo treinado, e mapas de modelos treinados com sementes diferentes concordam apenas moderadamente~\cite{arends2026}; o piso mede a variação do método, e não a do treino. BRE (20 mapas) e STE (14) têm posições menos estáveis. O CPSC2018 não anota as regiões diagnósticas, de modo que a plausibilidade clínica das explicações não foi avaliada, e ela pode discordar da fidelidade~\cite{storas2025}. Os limites das regiões são aproximados, e o teste de robustez cobriu apenas a \CS\ e a \CX. A soma entre batimentos apaga variações de um batimento para outro, como a irregularidade do RR na AF ou um batimento prematuro isolado na CAP e na CVP. Por fim, o ruído do GradientSHAP não foi ajustado para o ECG.
